\introduction

El sector energético cubano~\citep{ministerio2023} enfrentaba retos significativos en la operación estable de su sistema eléctrico nacional. Estos retos se caracterizaban por una capacidad de generación que en ocasiones no satisfacía la demanda máxima, situación conocida como déficit de generación. Este escenario obligaba al Despacho Nacional de Carga (DNC) a implementar medidas de gestión de la demanda~\citep{union2023}, entre las que se destacaba la rotación controlada de circuitos de distribución.

La rotación de circuitos, esencial para equilibrar la red y evitar colapsos, implicaba la desconexión temporal y secuencial de circuitos no críticos, distribuidos geográficamente, para reducir la carga total del sistema durante períodos de máxima tensión. Esta tarea, compleja y de alta responsabilidad, recaía directamente en los especialistas del Despacho Provincial de Carga, quienes debían considerar un conjunto multifactorial de variables en tiempo real: consumo horario por circuito (\texttt{ap\_curvas}), estado de averías (\texttt{ap\_apagones}), circuitos bajo aseguramiento (\texttt{ap\_Aseguramientos}) y restricciones operativas de cada bloque.

Históricamente, la generación del orden de rotación de circuitos se realizaba de forma manual. Los especialistas analizaban planillas con las cargas históricas y pronosticadas de cada circuito (tabla \texttt{ap\_curvas}), los megawatt-hora (MWh) totales, y las solicitudes específicas del DNC~\citep{union2023}. Además, cruzaban esta información con el estado operativo de la red: circuitos averiados (\texttt{ap\_apagones}), bajo aseguramiento temporal (\texttt{ap\_Aseguramientos}), o afectados por mantenimientos planificados.

Esta metodología, aunque funcional, presentaba limitaciones inherentes. Era un proceso lento, propenso a errores humanos dada la gran cantidad de datos y su dinamismo~\citep{pressman2020}, y difícilmente reproducible bajo los mismos criterios por diferentes operadores. La subjetividad en la priorización y la falta de capacidad para simular rápidamente múltiples escenarios ante un cambio en las condiciones~\citep{zhou2023} eran factores que podían impactar negativamente en la eficiencia de la medida y, en última instancia, en la calidad del servicio brindado a la población.

Esta investigación se desarrolló en el contexto del Despacho Provincial de Carga, unidad operativa clave en la ejecución de las directivas nacionales. Se contaba con una infraestructura de datos que incluía bases de datos en SQL Server 2008~\citep{jones2019}, las cuales almacenaban información histórica de cargas y eventos de la red en tablas como \texttt{ap\_circuitos}, \texttt{ap\_curvas}, \texttt{ap\_apagones} y \texttt{ap\_Aseguramientos}. Sin embargo, no existía una plataforma integrada de arquitectura moderna que automatizara el núcleo del proceso de decisión. La brecha identificada era la falta de una herramienta tecnológica que, aprovechando los datos existentes y los criterios expertos, centralizara el análisis y generara propuestas de órdenes de rotación de manera rápida y consistente, accesible desde cualquier nodo de la red operativa. La solución propuesta implementa una arquitectura de microservicios con API Gateway y comunicación NATS, desacoplada de la base de datos legacy.

La rotación de circuitos eléctricos en el Despacho Provincial de Carga Cienfuegos, ante situaciones de déficit de generación en la red de distribución del sistema electroenergético cubano, se realizaba manualmente sin un análisis integral de las múltiples variables que intervenían en el proceso en tiempo real. Por tanto, surgió la siguiente pregunta de investigación: ¿cómo automatizar el proceso de rotación de circuitos eléctricos en el Despacho Provincial de Carga Cienfuegos mediante un sistema de microservicios con API Gateway que analice variables en tiempo real (consumos, apagones, aseguramientos), para generar órdenes de desconexión y reconexión óptimas que reduzcan el tiempo de respuesta y eliminen los errores humanos en el cálculo del déficit?

Se planteó la hipótesis de que si se implementaba un sistema de microservicios (API Gateway con NestJS, circuitos-ms, rotaciones-ms) que analizara variables en tiempo real (cargas horarias, apagones, aseguramientos), entonces se reduciría el tiempo de respuesta operativa de 30–45 minutos a menos de 5 segundos y se eliminarían los errores humanos en el cálculo del déficit energético, optimizando así la generación de órdenes de rotación de circuitos en el Despacho Provincial de Carga.

El objetivo general de este trabajo fue diseñar e implementar un sistema de microservicios que generara las órdenes de circuitos a rotar ante déficits de generación, optimizando el proceso de decisión del Despacho Provincial. Para ello, se propuso desarrollar una Single Page Application (SPA) utilizando Next.js y React para la capa de presentación, estilizada con Tailwind CSS para garantizar una interfaz de usuario (UI) intuitiva y reactiva. El backend se implementó como una arquitectura de microservicios con NestJS, compuesta por un API Gateway que actúa como punto de entrada unificado y microservicios especializados: \textbf{circuitos-ms} (gestión de circuitos y consumos), \textbf{rotaciones-ms} (algoritmo de rotación FIFO por tiempo de estado) y \textbf{auth-ms} (autenticación JWT con Redis). Los microservicios se comunican internamente mediante \textbf{NATS} y consumen las tablas maestras de SIGERE (\texttt{ap\_circuitos}, \texttt{ap\_curvas}, \texttt{ap\_apagones}) y PSFV (\texttt{ap\_Aseguramientos}) en SQL Server 2008. El sistema permitía la supervisión, ajuste y validación final por parte del especialista a través de una interfaz dinámica que gestionaba el estado de los circuitos mediante Hooks y Context API, transformando su rol de ejecutor manual a supervisor de propuestas generadas automáticamente.

Como objetivos específicos se propusieron:
\begin{itemize}
    \item Establecer los fundamentos teórico-metodológicos del proceso de desarrollo del sistema de microservicios que sustentan la investigación.
    \item Desarrollar los fundamentos conceptuales y metodológicos asociados a la automatización de procesos en sistemas de gestión energética, relacionados con la operación del sistema electroenergético cubano en escenarios de déficit de generación.
    \item Realizar el diagnóstico integral del proceso manual actual de rotación de circuitos en el Despacho Provincial de Carga Cienfuegos.
    \item Establecer los fundamentos tecnológicos que sustentan el desarrollo del sistema de microservicios de rotación de circuitos, integrando las consideraciones teóricas y metodológicas para la selección de la plataforma tecnológica (NestJS, NATS, Next.js, React, SQL Server 2008).
    \item Caracterizar la solución tecnológica desarrollada: arquitectura de microservicios, API Gateway, circuitos-ms, rotaciones-ms, auth-ms y frontend SPA.
    \item Sistematizar los resultados del proceso de Ingeniería de Requisitos aplicado al sistema de microservicios de rotación de circuitos.
    \item Diseñar la arquitectura de datos para el almacenamiento, procesamiento y transmisión de datos, y la implementación de la interfaz gráfica basada en tecnologías web.
    \item Verificar y validar los resultados obtenidos con la implementación del sistema de microservicios.
\end{itemize}

La justificación de esta investigación radicaba en su potencial impacto práctico y económico. Al reducir el tiempo de respuesta ante una alerta de déficit, se aumentaba la resiliencia operativa. Además, la transición a una arquitectura de microservicios con API Gateway y NATS eliminaba las barreras de despliegue y mantenimiento de software de escritorio tradicional, permitiendo escalabilidad horizontal y despliegue centralizado. Metodológicamente, el proyecto se enmarcó dentro del desarrollo de software aplicado, empleando un enfoque ágil basado en la Programación Extrema (XP)~\citep{beck2000}, lo que permitió una adaptación continua a los requerimientos del cliente y la entrega de una herramienta escalable y moderna para la infraestructura crítica eléctrica.